\chapter{Praxislabor: zehn Experimente}
Führe jedes Experiment in einem neuen Chat aus und notiere: Auftrag, verwendete Werkzeuge, Rückfragen, Ergebnis, Fehler und Verbesserung.
\section{1. Ein Prompt, drei Qualitätsstufen}
Bitte zuerst nur um eine Erklärung. Füge dann Zielgruppe und Format hinzu. Formuliere schließlich Erfolgskriterien. Vergleiche die Ergebnisse.
\section{2. Eigener Command}
Definiere \texttt{/brief}: Ziel, Ausgangslage, drei Kernaussagen, Risiken, nächste Schritte. Teste ihn an zwei Themen und verbessere die Definition.
\section{3. Modifier kombinieren}
Teste \texttt{/vergleich +kurz}, danach \texttt{+tief +kritisch +tabelle}. Beobachte, welche Modifier miteinander konkurrieren.
\section{4. Aktualität erzwingen}
Stelle eine zeitabhängige Frage. Verlange Datum, Quellen und Kennzeichnung unsicherer Aussagen. Öffne mindestens zwei Quellen.
\section{5. Deep Research}
Wähle eine Frage mit mehreren Perspektiven. Begrenze Zeitraum und Quellentypen. Prüfe den Forschungsplan vor dem Start.
\section{6. Datei als Datenquelle}
Lade ein eigenes, unkritisches Dokument hoch. Lass Aussagen extrahieren und fordere Seiten- oder Abschnittsnachweise. Suche eine erfundene Behauptung.
\section{7. Workflow mit Zwischenschritt}
Lass aus Notizen erst eine Gliederung erstellen. Genehmige sie, bevor der Text entsteht. Vergleiche das Resultat mit einer Ein-Schritt-Anfrage.
\section{8. Read--Draft--Act}
Formuliere eine Aufgabe dreimal: nur lesen, Entwurf erstellen, Aktion vorbereiten. Markiere, an welcher Stelle eine Bestätigung erforderlich ist.
\section{9. Fehler provozieren}
Gib absichtlich widersprüchliche Randbedingungen. Prüfe, ob Rückfragen kommen oder Annahmen transparent gemacht werden.
\section{10. Eigener wiederverwendbarer Workflow}
Wähle eine monatliche Aufgabe. Definiere Eingangsdaten, Schritte, Qualitätsprüfung, Freigabepunkte und Ausgabe. Teste mit einem realen Beispiel.
